\documentclass[aspectratio=169]{beamer}

% Compilar con XeLaTeX
\usepackage{fontspec}
\usepackage{xcolor}
\usepackage{graphicx}
\usepackage{listings}

% Branding Académico (Obligatorio)
\definecolor{AzulMarino}{HTML}{1A365D}
\definecolor{GrisPizarra}{HTML}{4A5568}
\definecolor{VerdeAzulado}{HTML}{319795}
\definecolor{FondoBlanco}{HTML}{F7FAFC}

% Configuracion de Beamer
\setbeamercolor{background canvas}{bg=FondoBlanco}
\setbeamercolor{title}{fg=AzulMarino}
\setbeamercolor{frametitle}{fg=AzulMarino}
\setbeamercolor{normal text}{fg=GrisPizarra}

% Fuentes
\setmainfont{Inter}
\setsansfont{Montserrat}
\setmonofont{JetBrains Mono}

% Configuración de listings para código Python
\lstset{
    language=Python,
    basicstyle=\ttfamily\small,
    keywordstyle=\color{AzulMarino}\bfseries,
    stringstyle=\color{VerdeAzulado},
    commentstyle=\color{GrisPizarra}\itshape,
    showstringspaces=false,
    breaklines=true,
    frame=single,
    rulecolor=\color{GrisPizarra}
}

\title{Práctica: Modelos Supervisados Clásicos}
\subtitle{Scikit-Learn, Pandas y Visualización}
\author{Machine Learning para Ciberseguridad}
\date{Semana 3}

\begin{document}

\begin{frame}
    \titlepage
\end{frame}

\begin{frame}[fragile]{Exploración del Dataset (Pandas)}
    Vamos a usar \texttt{pandas} para cargar e inspeccionar la data:
    \begin{lstlisting}
import pandas as pd

# Cargar dataset de URLs sospechosas
df = pd.read_csv("phishing_dataset.csv")

# Explorar las primeras filas
print(df.head())
    \end{lstlisting}
\end{frame}

\begin{frame}[fragile]{Entrenamiento: k-NN en Scikit-Learn}
    La librería \texttt{scikit-learn} tiene una interfaz estándar: \texttt{fit()} para entrenar y \texttt{predict()} para clasificar.
    \begin{lstlisting}
from sklearn.neighbors import KNeighborsClassifier

# Instanciar el modelo (k=5 vecinos)
knn = KNeighborsClassifier(n_neighbors=5, metric='euclidean')

# Entrenar el modelo con los datos
knn.fit(X_train, y_train)

# Clasificar nuevas paginas web
predicciones = knn.predict(X_test)
    \end{lstlisting}
\end{frame}

\begin{frame}[fragile]{Evaluación: La Matriz de Confusión}
    \begin{lstlisting}
from sklearn.metrics import confusion_matrix, ConfusionMatrixDisplay
import matplotlib.pyplot as plt

# Generar la matriz cruzando realidad vs prediccion
cm = confusion_matrix(y_test, predicciones)

# Ploteo visual
disp = ConfusionMatrixDisplay(cm, display_labels=['Phishing', 'Legitimo'])
disp.plot(cmap='Blues')
plt.show()
    \end{lstlisting}
\end{frame}

\end{document}
